\documentclass[10pt,a4paper]{amsart}
\usepackage[margin=19mm]{geometry}
\usepackage{amsmath,amssymb,mathtools}
\usepackage[scr=rsfs,cal=euler]{mathalfa}
\usepackage{xfrac}
\usepackage[unicode,hidelinks,bookmarks=false]{hyperref}
\newcommand{\ie}{i.e.\ }
\newcommand{\cf}{cf.\ }
\newcommand{\resp}{respectively\ }
\newcommand{\Subsection}{\subsection}
\providecommand{\nobreakdash}{\nobreak\hbox{-}}
\setlength{\emergencystretch}{2em}
\multlinegap0pt
\usepackage{bm}
\usepackage{bbm}
\usepackage{dsfont}
\usepackage{xspace}
\usepackage{cancel}
\usepackage{mathtools}
\usepackage{amsthm}
\makeatletter
\@ifundefined{thm}{\newtheorem{thm}{Thm}}{}
\@ifundefined{lem}{\newtheorem{lem}[thm]{Lemma}}{}
\makeatother
\makeatletter
\expandafter\ifx\csname AMS\endcsname\relax
\newenvironment{AMS}{}{}
\fi
\newcommand{\F}{\mathbb{F}}
\DeclareMathOperator{\diag}{diag}
\expandafter\ifx\csname keywords\endcsname\relax
\newenvironment{keywords}{}{}
\fi
\DeclareMathOperator{\rank}{rank}
\makeatother
\title[On gcd-graphs over matrix rings]{On gcd-graphs over matrix rings}
\author{Dung Nguyen, Tung T. Nguyen}
\date{Source: arXiv:2606.00836v1 (CC BY 4.0)}
\begin{document}
\maketitle

\noindent\emph{Source and licence.} On gcd-graphs over matrix rings. Version arXiv:2606.00836v1 (CC BY 4.0), distributed under the Creative Commons Attribution 4.0 International licence, as recorded in the arXiv licence metadata for this exact version and shown on the abstract page https://arxiv.org/abs/2606.00836v1. Full rights notice: licenses/arxiv-2606.00836-CC-BY-4.0.txt.\par\smallskip

\noindent\textit{Editorial scope.} Counted unit: the Lem whose source label is \texttt{lem:r<m} in the article (source0.tex lines 222--224, source label \texttt{lem:r<m}), delivered together with the article's own complete original proof of it (source0.tex lines 225--304, one \texttt{proof} environment of 80 lines). No mathematics is written, repaired or re-derived: statement and proof are the article's own text line for line. Required context retained uncounted: none. Every definition, display and label of those lines is retained unchanged. Omitted from this package and not redistributed: 1--221, 305--820. Nothing inside a retained excerpt is omitted or altered: every retained body is delivered exactly as the article's lines are written, so no omission receipt of any kind accompanies this package. Citations: the retained excerpts carry no citation command at all, so no bibliography entry of the article is transcribed and no reference is redistributed. Reference adaptation: none was needed and none was made; every reference inside the delivered unit resolves inside it, so no unresolved reference or citation is produced. Printed slips kept literal: no printed word, symbol, hypothesis or number of the counted statement or of its proof was altered. Layout and class: the article's own class is \texttt{amsart} with packages including geometry, amsmath, amssymb, amsfonts, float, hyperref, cleveref, graphicx, xcolor, multicol, tabularx, mathpazo; the wrapper is the lane's pinned \texttt{amsart} editorial wrapper at 10pt on A4, which restates the theorem-like environments the retained text needs and adds the article's own macro definitions that the retained text needs (\texttt{\textbackslash F}, \texttt{\textbackslash diag}, \texttt{\textbackslash rank}), with extra wrapper packages (bm, bbm, dsfont, xspace, cancel, mathtools). The delivered numbering is the unit's own. Assets: no retained span contains a figure environment, a graphics-inclusion command, a PDF-page inclusion, a table or a TikZ picture; no image file is copied into this package, the package declares no asset dependency and no image file is redistributed with it. Licence: version arXiv:2606.00836v1 of this article is distributed under the Creative Commons Attribution 4.0 International licence, as recorded in the arXiv licence metadata for this exact version and shown on the abstract page https://arxiv.org/abs/2606.00836v1. A full rights notice is delivered at licenses/arxiv-2606.00836-CC-BY-4.0.txt. Characters: every retained excerpt is pure ASCII, so the delivered engine, XeLaTeX with its default Latin Modern text font, reports no missing character.
\begin{lem} \label{lem:r<m}
    If $m-k \leq r \leq m$, any $n \times n$ matrix with rank $r$ is a sum of an $n \times n$ matrix with rank $m$ and an $n \times n$ matrix with rank $k$.
\end{lem}

\begin{proof}
    First, let us suppose that $F$ has more than 2 elements. Then, we may choose an $h \in F \setminus\{0,1\}$ and write $I_n(r)$ as a sum of 2 matrices
    \begin{align*}
        I_n(r)
        &= \diag(\underbrace{1,\ldots,1}_{m-k},\underbrace{1-h,\ldots,1-h}_{r-(m-k)},\underbrace{1,\ldots,1}_{m-r},\underbrace{0,\ldots,0}_{n-m})\\
        &+ \diag(\underbrace{0,\ldots,0}_{m-k},\underbrace{h,\ldots,h}_{r-(m-k)},\underbrace{-1,\ldots,-1}_{m-r},\underbrace{0,\ldots,0}_{n-m}).
    \end{align*}
    Let $X_m, X_k$ be the two diagonal matrices in the equation above. Then, $X_m$ has $m$ non-zero rows and thus has rank $m$ while $X_k$ has $k$ non-zero rows and thus has rank $k$.
    If $F$ is the field with two elements, we further divide the condition into two cases, $m>k$ or $m=k$.
    \paragraph{Case 1: $m>k$}
    Let $U_r$ be an $r\times r$ upper triangular matrix with ones on the main diagonal. We set exactly $r - (m-k)$ ones on the superdiagonal of $U_r$, and the rest to zeros. Because $m > k$, we are guaranteed that $r - (m-k) \le r-1$, ensuring there are enough positions on the superdiagonal. We can see that  that $U_r$ has rank $r$ and it has the following explicit description 
    \[U_r = I_r +\sum_{i=1}^{r - (m-k)} E_{i,(i+1)}.\]
    We then construct the submatrices $X_m,X_k$ as follows:
    \[
    X_m =
    \begin{bmatrix}
        U_r & 0\\
        0 & I_{m-r}
    \end{bmatrix},
    \]
    \[
    X_k = I_m(r) - X_m =
    \begin{bmatrix}
        I_r-U_r & 0\\
        0 & -I_{m-r}
    \end{bmatrix}.
    \]
    Utilizing the properties of block diagonal matrices,
    \[\rank(X_m) = \rank(U_r) + \rank(I_{m-r}) = r+m-r=m.\]
    The block $I_r - U_r$ is  upper triangular with exactly $r - (m-k)$ non-zero entries on its superdiagonal, giving it a rank of exactly $r - (m-k)$. 
    \[\rank(X_k) = \rank(I_r - U_r) + \rank(-I_{m-r}) = r - (m-k)+m-r=k.\]
    Adding $X_m,X_k$ we get $I_r$ as the only non-zero block of the sum,
    \[
    X_m + X_k =
    \begin{bmatrix}
        U_r + I_r - U_r & 0\\
        0 & I_{m-r} - I_{m-r}
    \end{bmatrix}
    =
    \begin{bmatrix}
        I_r& 0\\
        0 & 0
    \end{bmatrix}.
    \]
    Both $X_m$ and $X_k$ are then padded with zeros to reach the dimension $n\times n$.

    \paragraph{Case 2: $m=k$}
    We consider 3 more subcases, either $r<m$, $r=m \geq 2$, or $r=m=1$.
    \subparagraph{Case 2.1: $r<m$}
    Let $X_m$ be the following $m\times m$ cyclic shift matrix:
    \[X_m = \left[ \sum_{i=1}^{m-1} E_{i,(i+1)} \right] + E_{m,1}.\]
    Evaluating the determinant of $X_m$ by expanding along the first column isolates the single non-zero entry coming from $E_{m,1}$. The corresponding minor is $I_{m-1}$, yielding a non-zero determinant. Thus $X_m$ is invertible and has full rank $m$. Then, let $X_k = I_m(r) - X_m$. Since $r<m$, the ones from $I_m(r)$ do not appear on the $m$-th row of $X_k$. Evaluating the determinant of $X_k$ by expanding along the $m$-th (bottom) row, the only non-zero entry comes from $E_{m,1}$. The corresponding minor is an upper triangular matrix with ones all across its main diagonal, yielding a non-zero determinant. Thus $X_k$ is also invertible and has full rank $m$. $X_m$ and $X_k$ add up to $I_m(r)$.
    % \tung{Why $X_k$ has rank $m$? Where do we use the condition $r<m$?}
    \subparagraph{Case 2.2: $r=m \geq 2$}
    We make one small adjustment to $X_m$:
    \[X_m = \left[ \sum_{i=1}^{m-1} E_{i,(i+1)} \right] + E_{m,1} + E_{m,2}.\]
    Similarly, through the calculation of the determinant with the first column, we can show that $X_m$ has full rank $m$. Then, let $X_k = I_m - X_m$. The matrix $X_k$ has ones on its main diagonal, superdiagonal, and at positions $(m,1)$ and $(m,2)$. Because the first row of $X_k$ contains ones exactly at $(1,1)$ and $(1,2)$ with zeros elsewhere, we can apply the elementary row operation of adding the first row to the $m$-th row. Over $\F_2$, this turn the entries at $(m,1)$ and $(m,2)$ to zero, leaving the $m$-th row with only a single one located on the main diagonal at $(m,m)$. This transforms $X_k$ into an upper triangular matrix with ones completely across its main diagonal, thus has full rank $m$. Because elementary row operations preserve the rank, $X_k$ has full rank $m$. $X_m$ and $X_k$ add up to $I_m$, and thus $I_m(r)$ since $r=m$.
    \subparagraph{Case 2.3: $r=m=1$}
    We default to the following rank 1 matrices $X_m,X_k$:
    \[
    X_m = 
    \begin{bmatrix}
        1 & 1\\
        0 & 0
    \end{bmatrix},
    X_k =
    \begin{bmatrix}
        0 & 1\\
        0 & 0
    \end{bmatrix},
    \]
    \[
    X_m+X_k =
    \begin{bmatrix}
        1 & 0\\
        0 & 0
    \end{bmatrix}.
    \]
    Both $X_m$ and $X_k$ are then padded with zeros to reach the dimension $n\times n$.
\end{proof}

\end{document}
